\documentclass[10pt,a4paper]{article}

% ============================================================
% A4 - PLAN DE GESTION DE DATOS ROUTTECH
% Documento autocontenido y compatible con pdfLaTeX en Overleaf.
% No requiere archivos .sty, configuracion.tex ni subcarpetas.
% ============================================================

\usepackage[utf8]{inputenc}
\usepackage[T1]{fontenc}
\usepackage[spanish,es-nodecimaldot]{babel}
\usepackage{tgheros}
\renewcommand{\familydefault}{\sfdefault}
\usepackage[a4paper,top=1.85cm,bottom=1.8cm,left=1.9cm,right=1.9cm,headheight=20pt]{geometry}
\usepackage{xcolor}
\usepackage{tabularx,array,booktabs,colortbl,longtable}
\usepackage{enumitem}
\usepackage[most]{tcolorbox}
\usepackage{fancyhdr}
\usepackage{lastpage}
\usepackage{microtype}
\usepackage{setspace}
\usepackage{titlesec}
\usepackage{hyperref}
\usepackage{xurl}
\usepackage{tikz}
\usepackage{ragged2e}
\usepackage{amssymb}

% ---------- PALETA PROPIA ROUTTECH ----------
\definecolor{RTForest}{HTML}{25463B}
\definecolor{RTTerracotta}{HTML}{B45F45}
\definecolor{RTGold}{HTML}{D5A64A}
\definecolor{RTCream}{HTML}{F7F2E8}
\definecolor{RTLight}{HTML}{EEF3F0}
\definecolor{RTGray}{HTML}{4D5652}
\definecolor{RTRed}{HTML}{A23B3B}

% ---------- DATOS DEL DOCUMENTO ----------
\newcommand{\Proyecto}{RoutTech: Sistema de Movilidad Colaborativa para la Comunidad Universitaria de la UTEQ}
\newcommand{\Asignatura}{ISR-401 -- Ingeniería de Requisitos}
\newcommand{\Periodo}{2026--2027 PPA}
\newcommand{\CategoriaEtica}{Categoría B -- Datos personales}
\newcommand{\Repositorio}{https://github.com/AnderJM3/RoutTech_ISR401}
\newcommand{\FechaDocumento}{29 de julio de 2026}
\newcommand{\VersionDocumento}{2.0}
\newcommand{\Paralelo}{Cuarto nivel -- Software B}
\newcommand{\EstadoDocumento}{Versión 2.0 para revisión ética y firma}

\newcommand{\Docente}{Ing. Gleiston Guerrero Ulloa, PhD}
\newcommand{\CorreoDocente}{gguerrero@uteq.edu.ec}

\newcommand{\Anderson}{NARANJO FLORES ANDERSON JEAMPIERE}
\newcommand{\CedulaAnderson}{1207412659}
\newcommand{\MatriculaAnderson}{617251}
\newcommand{\CorreoAnderson}{anaranjof2@uteq.edu.ec}

\newcommand{\Mayra}{CORDOVA CARREÑO MAYRA LUCILA}
\newcommand{\CedulaMayra}{0969483943}
\newcommand{\MatriculaMayra}{617891}
\newcommand{\CorreoMayra}{mcordovac2@uteq.edu.ec}

% ---------- TABLAS ----------
\newcolumntype{Y}{>{\RaggedRight\arraybackslash}X}
\newcolumntype{J}{>{\justifying\arraybackslash}X}
\newcolumntype{C}[1]{>{\centering\arraybackslash}p{#1}}
\newcolumntype{L}[1]{>{\RaggedRight\arraybackslash}p{#1}}
\newcolumntype{P}[1]{>{\justifying\arraybackslash}p{#1}}
\renewcommand{\arraystretch}{1.20}
\setlength{\tabcolsep}{4.2pt}

% ---------- ENCABEZADO Y PIE ----------
\pagestyle{fancy}
\fancyhf{}
\fancyhead[L]{\small\textbf{UTEQ | FCC | ISR-401}}
\fancyhead[R]{\small\textbf{A4 -- Plan de gestión de datos}}
\fancyfoot[L]{\small Paquete ético RoutTech}
\fancyfoot[R]{\small Página \thepage\ de \pageref{LastPage}}
\renewcommand{\headrulewidth}{0.8pt}
\renewcommand{\footrulewidth}{0.4pt}
\renewcommand{\headrule}{\hbox to\headwidth{\color{RTTerracotta}\leaders\hrule height \headrulewidth\hfill}}
\renewcommand{\footrule}{\hbox to\headwidth{\color{RTGold}\leaders\hrule height \footrulewidth\hfill}}

% ---------- TÍTULOS ----------
\titleformat{\section}
  {\large\bfseries\color{RTForest}}
  {\thesection.}{0.45em}{}
  [\vspace{1pt}\color{RTGold}\titlerule]
\titleformat{\subsection}
  {\normalsize\bfseries\color{RTTerracotta}}
  {\thesubsection}{0.45em}{}
\titlespacing*{\section}{0pt}{0.75em}{0.35em}
\titlespacing*{\subsection}{0pt}{0.55em}{0.25em}

% ---------- CAJAS ----------
\newtcolorbox{InfoBox}[1][]{
  enhanced,breakable,
  colback=RTLight,colframe=RTForest,
  boxrule=0.8pt,arc=2mm,
  left=2.6mm,right=2.6mm,top=1.6mm,bottom=1.6mm,
  title=#1,fonttitle=\bfseries,coltitle=white,
  colbacktitle=RTForest
}
\newtcolorbox{WarningBox}[1][]{
  enhanced,breakable,
  colback=RTCream,colframe=RTTerracotta,
  boxrule=0.8pt,arc=2mm,
  left=2.6mm,right=2.6mm,top=1.6mm,bottom=1.6mm,
  title=#1,fonttitle=\bfseries,coltitle=white,
  colbacktitle=RTTerracotta
}

\hypersetup{
  colorlinks=true,
  linkcolor=RTForest,
  urlcolor=RTTerracotta,
  pdftitle={A4 Plan de Gestion de Datos RoutTech},
  pdfauthor={Equipo RoutTech}
}
\urlstyle{same}
\setlength{\parindent}{1.1em}
\setlength{\parskip}{2.5pt}
\setstretch{1.10}
\setlist[itemize]{leftmargin=1.35em,itemsep=1pt,topsep=2pt}
\setlist[enumerate]{leftmargin=1.55em,itemsep=1pt,topsep=2pt}

\begin{document}
\justifying

% ======================= PORTADA =======================
\thispagestyle{empty}
\begin{tikzpicture}[remember picture,overlay]
  \fill[RTForest] (current page.north west) rectangle ([yshift=-4.6cm]current page.north east);
  \fill[RTTerracotta] ([yshift=-4.6cm]current page.north west) rectangle ([yshift=-4.95cm]current page.north east);
\end{tikzpicture}

\vspace*{0.65cm}
\begin{center}
{\color{white}\Large\bfseries UNIVERSIDAD TÉCNICA ESTATAL DE QUEVEDO}

\vspace{1.75cm}
{\large\bfseries\color{RTForest} Facultad de Ciencias de la Computación\\[4pt]}
{\normalsize\color{RTGray} Carrera de Ingeniería de Software}

\vspace{0.8cm}

{\Huge\bfseries\color{RTForest} PLAN DE GESTIÓN DE DATOS\\[8pt]}
{\Large\bfseries\color{RTTerracotta} A4 -- Data Management Plan}

\vspace{1.0cm}

\begin{tcolorbox}[width=0.90\textwidth,colback=RTCream,colframe=RTGold,boxrule=1pt,arc=2mm]
\centering
{\large\bfseries \Proyecto}\\[7pt]
\textbf{Clasificación ética:} \CategoriaEtica\\
\textbf{Asignatura:} \Asignatura\\
\textbf{Periodo académico:} \Periodo
\end{tcolorbox}

\vfill
\begin{tabular}{rl}
\textbf{Documento:} & A4 -- Plan de gestión de datos\\
\textbf{Versión:} & \VersionDocumento\\
\textbf{Estado:} & \EstadoDocumento\\
\textbf{Repositorio:} & \href{\Repositorio}{RoutTech\_ISR401}\\
\textbf{Fecha:} & \FechaDocumento
\end{tabular}

\vspace{0.75cm}
{\small Quevedo -- Los Ríos -- Ecuador}
\end{center}
\clearpage

% ======================= CONTENIDO =======================
\begin{WarningBox}[Principio de separación de datos]
El repositorio \href{\Repositorio}{RoutTech\_ISR401} se utilizará para documentación, instrumentos, matrices, código y materiales completamente anonimizados. Los consentimientos firmados, cédulas, grabaciones identificables y archivos crudos no se almacenarán ni publicarán en GitHub; se conservarán únicamente en carpetas locales cifradas bajo custodia del equipo investigador.
\end{WarningBox}

\section{Identificación y alcance del plan}

\rowcolors{2}{RTLight}{white}
\noindent\begin{tabularx}{\textwidth}{@{}L{4.45cm} J@{}}
\rowcolor{RTForest}
\textcolor{white}{\textbf{Campo}} & \textcolor{white}{\textbf{Información}}\\
Proyecto & \Proyecto.\\
Categoría ética & \CategoriaEtica.\\
Equipo custodio & \Anderson\ y \Mayra.\\
Docente responsable & \Docente\ -- \href{mailto:\CorreoDocente}{\CorreoDocente}.\\
Paralelo & \Paralelo.\\
Repositorio del proyecto & \url{\Repositorio}.\\
Vigencia del plan & Desde la aprobación interna del protocolo hasta la destrucción o anonimización definitiva de los datos.\\
Objetivo & Definir cómo se recopilarán, nombrarán, almacenarán, protegerán, respaldarán, utilizarán, compartirán, conservarán y eliminarán los datos del proyecto RoutTech.\\
\end{tabularx}

\section{Inventario y clasificación de los datos}

\small
\rowcolors{2}{RTCream}{white}
\noindent\begin{tabularx}{\textwidth}{@{}L{2.55cm} L{3.0cm} J L{2.45cm}@{}}
\rowcolor{RTTerracotta}
\textcolor{white}{\textbf{Tipo}} &
\textcolor{white}{\textbf{Formato}} &
\textcolor{white}{\textbf{Contenido}} &
\textcolor{white}{\textbf{Clasificación}}\\
Consentimientos & PDF o imagen digitalizada & Nombre, cédula, firma, fecha y autorizaciones del participante. & Restringido\\
Entrevistas & MP4, MP3/WAV, DOCX/TXT & Imagen, voz, respuestas, transcripción y metadatos. & Restringido hasta anonimizar\\
Cuestionario & XLSX/CSV y resumen PDF & Respuestas de Google Forms; no se publicarán identificadores directos. & Interno / anonimizado\\
Observación & DOCX/XLSX, JPG/PNG autorizados & Notas de campo y fotografías generales sin rostros, placas ni domicilios. & Interno / publicable tras revisión\\
Walkthrough & MP4, PDF, XLSX & Grabación autorizada, acta, observaciones y evaluación de requisitos o mockups. & Restringido hasta anonimizar\\
Artefactos técnicos & TEX, PDF, CSV, JSON, PNG/SVG, PUML & Requisitos, matrices, diagramas, mockups, scripts y código del MVP. & Publicable tras revisión\\
\end{tabularx}
\normalsize

\begin{InfoBox}[Volumen estimado]
Considerando dieciséis entrevistas, audios, transcripciones, documentos y sesiones de validación, se prevé un volumen aproximado de \textbf{20 a 40 GB}. Este rango es una estimación operativa y se actualizará si el tamaño real de los archivos supera lo previsto.
\end{InfoBox}

\section{Recolección y minimización}

La recolección se limitará a la información necesaria para elicitar y validar requisitos de RoutTech. Mayra Cordova dirigirá el trabajo de campo y verificará que exista consentimiento previo; Anderson Naranjo controlará la organización, integridad y trazabilidad de los archivos. Se aplicarán las siguientes reglas:

\begin{itemize}
  \item No solicitar contraseñas, credenciales institucionales reales, datos bancarios, placas reales ni direcciones domiciliarias exactas.
  \item No activar seguimiento de geolocalización en tiempo real durante el estudio académico.
  \item Registrar únicamente el perfil general del participante y los metadatos necesarios para la trazabilidad.
  \item Usar códigos seudónimos como \texttt{EV-ENT-001}, evitando nombres en archivos de trabajo y análisis.
  \item Recoger audio, video o fotografías solo cuando exista autorización expresa en el consentimiento.
  \item Utilizar datos ficticios o sintéticos en el MVP, demostraciones, pruebas y capturas públicas.
\end{itemize}

\section{Convención de nombres y metadatos}

\begin{InfoBox}[Formato de nombre]
\texttt{AAAA-MM-DD\_Tipo\_Codigo\_Perfil\_vN.ext}

Ejemplo: \texttt{2026-07-29\_Entrevista\_EV-ENT-001\_Estudiante\_v1.mp4}
\end{InfoBox}

Cada evidencia deberá registrar como mínimo: código, fecha, tipo de técnica, perfil general, responsable de captura, versión, formato, duración o tamaño, estado de consentimiento, estado de anonimización y ubicación de almacenamiento. Los cambios relevantes generarán una nueva versión; no se sobrescribirán archivos originales sin mantener una copia verificable.

\section{Almacenamiento, seguridad y acceso}

\rowcolors{2}{RTLight}{white}
\noindent\begin{tabularx}{\textwidth}{@{}L{3.25cm} J L{3.15cm}@{}}
\rowcolor{RTForest}
\textcolor{white}{\textbf{Zona}} &
\textcolor{white}{\textbf{Contenido y medida de protección}} &
\textcolor{white}{\textbf{Acceso}}\\
Repositorio GitHub & Documentos, código, diagramas, instrumentos y datos completamente anonimizados. No contendrá cédulas, firmas ni grabaciones identificables. & Público según configuración del repositorio\\
Carpeta local cifrada principal & Datos crudos, consentimientos, videos, audios y transcripciones no anonimizadas. Cifrado mediante contenedor o archivo AES-256 protegido con contraseña robusta. & Anderson y Mayra\\
Respaldo cifrado & Copia sincronizada semanalmente en el segundo equipo autorizado o unidad externa cifrada, separada de la copia principal. & Anderson y Mayra\\
Entrega al docente & Documentos requeridos para revisión ética o académica, compartidos de forma directa y sin publicar enlaces abiertos. & Docente responsable\\
\end{tabularx}

Las contraseñas de cifrado no se almacenarán dentro del repositorio ni en archivos de texto junto a los datos. Los dispositivos deberán contar con bloqueo de sesión, sistema actualizado y protección antimalware. El intercambio de archivos identificables se evitará; cuando sea indispensable, se realizará mediante un canal institucional o archivo cifrado cuya clave se comunique por un medio diferente.

\section{Copias de seguridad, integridad y control de versiones}

\begin{enumerate}
  \item Se mantendrá una copia principal y una copia de respaldo cifrada, actualizada al menos una vez por semana durante el trabajo de campo.
  \item Se comprobará mensualmente que los archivos críticos puedan abrirse y que el respaldo no esté corrupto.
  \item Los documentos técnicos utilizarán control de versiones en Git; los datos crudos no se incorporarán al historial del repositorio.
  \item Las modificaciones de transcripciones y bases de datos producirán versiones numeradas, conservando el archivo original de solo lectura.
  \item Cuando sea necesario, se generarán sumas de verificación SHA-256 para comprobar la integridad de entregables y conjuntos anonimizados.
\end{enumerate}

\section{Anonimización y preparación para análisis}

La anonimización se realizará antes de integrar datos en informes o repositorios públicos. Se eliminarán nombres, cédulas, correos, teléfonos, firmas, rostros, voces no autorizadas, placas, domicilios, enlaces personales y combinaciones que permitan reidentificación. Las ubicaciones se generalizarán a campus, sector, barrio o parroquia. Las citas textuales se revisarán para retirar referencias personales o contextuales innecesarias.

Las entrevistas se codificarán en Microsoft Excel mediante categorías temáticas. Las respuestas del cuestionario de Google Forms se exportarán para calcular frecuencias y porcentajes, evitando publicar respuestas individuales que contengan información identificable. Los resultados se presentarán de forma agregada.

\section{Retención, eliminación y cierre}

\rowcolors{2}{RTCream}{white}
\noindent\begin{tabularx}{\textwidth}{@{}L{4.0cm} J L{2.75cm}@{}}
\rowcolor{RTTerracotta}
\textcolor{white}{\textbf{Elemento}} &
\textcolor{white}{\textbf{Política}} &
\textcolor{white}{\textbf{Plazo}}\\
Datos identificables y grabaciones & Conservación cifrada para verificación académica; eliminación segura al concluir el plazo o antes si dejan de ser necesarios. & Hasta 24 meses tras el cierre\\
Consentimientos firmados & Custodia restringida para demostrar autorización y atender solicitudes del participante. & Hasta 24 meses tras el cierre\\
Datos anonimizados y agregados & Podrán mantenerse como parte de la documentación académica y científica, siempre que no permitan reidentificación. & Conservación académica\\
Código, diagramas y documentos & Permanecerán en el repositorio del proyecto conforme a su licencia y a las decisiones del equipo y la UTEQ. & Indefinido mientras sean pertinentes\\
\end{tabularx}

La eliminación se realizará mediante borrado seguro de archivos y copias de respaldo. El equipo registrará la fecha, el tipo de datos eliminados, el responsable y el método aplicado. Si un dispositivo se reasigna o desecha, se eliminarán previamente los contenedores cifrados y sus copias.

\section{Publicación, reutilización y licencias}

Solo podrán publicarse instrumentos, requisitos, matrices, scripts, diagramas, datos agregados y transcripciones plenamente anonimizadas. No se publicarán consentimientos, cédulas, firmas, grabaciones identificables ni bases de datos crudas. Antes de cualquier depósito en GitHub, OSF o Zenodo, Anderson y Mayra efectuarán una revisión conjunta de anonimización.

Los materiales elaborados por el equipo podrán utilizar una licencia académica o Creative Commons cuando corresponda; dicha licencia no se extenderá a información de participantes ni a contenido institucional cuya publicación no haya sido autorizada.

\section{Derechos de los participantes}

Las personas participantes podrán solicitar información sobre el tratamiento de sus datos, corrección de datos identificables, retiro de su participación o eliminación de archivos vinculables mientras todavía sea posible identificarlos. Las solicitudes se dirigirán a \href{mailto:\CorreoAnderson}{\CorreoAnderson}, \href{mailto:\CorreoMayra}{\CorreoMayra} o al docente responsable en \href{mailto:\CorreoDocente}{\CorreoDocente}. Una vez que la información haya sido anonimizada e integrada en resultados agregados, podría no ser técnicamente posible separar una contribución individual.

\section{Gestión de incidentes}

Ante pérdida, acceso no autorizado, publicación accidental o alteración de datos, el integrante que detecte el incidente deberá:

\begin{enumerate}
  \item detener el acceso o difusión y preservar la evidencia del incidente;
  \item informar al docente responsable y al otro integrante del equipo en un plazo interno máximo de 72 horas;
  \item identificar los archivos, participantes y sistemas afectados;
  \item cambiar credenciales, revocar enlaces o eliminar copias expuestas, según corresponda;
  \item documentar causas, acciones correctivas y medidas para evitar recurrencia;
  \item contactar a los participantes afectados cuando la evaluación ética o institucional así lo requiera.
\end{enumerate}

\section{Responsabilidades}

\small
\rowcolors{2}{RTLight}{white}
\noindent\begin{tabularx}{\textwidth}{@{}L{3.9cm} L{3.45cm} J@{}}
\rowcolor{RTForest}
\textcolor{white}{\textbf{Responsable}} &
\textcolor{white}{\textbf{Rol}} &
\textcolor{white}{\textbf{Funciones de gestión de datos}}\\
\Anderson & Líder y administrador documental & Estructura de archivos, control de versiones, revisión de anonimización, integridad, respaldo, GitHub, LaTeX y cierre del plan.\\
\Mayra & Responsable de trabajo de campo y requisitos & Consentimiento, captura de evidencias, asignación de códigos, metadatos, transcripción, codificación temática y reporte inmediato de incidentes.\\
\Docente & Supervisor académico & Revisión ética y metodológica, orientación ante incidentes y autorización de entregables institucionales.\\
\end{tabularx}
\normalsize

\section{Aprobación del plan}

\begin{InfoBox}[Declaración]
Quienes suscriben declaran conocer este Plan de Gestión de Datos y se comprometen a aplicarlo durante la recolección, análisis, almacenamiento, publicación y cierre del proyecto RoutTech.
\end{InfoBox}

\vspace{0.8cm}
\noindent\begin{tabularx}{\textwidth}{@{}Y Y@{}}
\centering\rule{6.1cm}{0.4pt} & \centering\arraybackslash\rule{6.1cm}{0.4pt}\\
\centering \Anderson & \centering\arraybackslash \Mayra\\
\centering Líder del equipo & \centering\arraybackslash Responsable de trabajo de campo\\[4pt]
\centering C.I. \CedulaAnderson & \centering\arraybackslash C.I. \CedulaMayra\\[4pt]
\centering Fecha: \rule{2.7cm}{0.4pt} & \centering\arraybackslash Fecha: \rule{2.7cm}{0.4pt}
\end{tabularx}

\vspace{1.1cm}
\begin{center}
\rule{7.0cm}{0.4pt}\\
\Docente\\
Docente responsable\\[4pt]
Fecha: \rule{3.0cm}{0.4pt}
\end{center}

\end{document}
